\chapter{Závěr}
\label{chap:conclusion}

Cílem mé odborné praxe ve společnosti AVE Soft s.r.o. bylo zapojení se do reálného vývojového procesu a realizace konkrétních softwarových úkolů v rámci systému Evolio. Hlavním výstupem byl návrh a implementace webového rozhraní pro modul Importér, který nahrazuje zastaralé desktopové řešení.

Během praxe se podařilo vytvořit moderní Single Page Aplikaci v technologii Angular, která komunikuje s nově vytvořeným backendovým doplňkem. Ačkoliv je modul momentálně ve fázi finalizace a testování, již nyní nabízí intuitivnější uživatelské rozhraní a zjednodušuje proces importu dat. Vedlejší úkoly, jako implementace AI chatu a tvorba automatizovaných testů, byly úspěšně dokončeny a nasazeny do provozu, kde přinášejí reálnou hodnotu jak firmě, tak jejím klientům.

Tato zkušenost pro mě byla klíčová pro pochopení fungování softwarové firmy. Osvojil jsem si moderní technologie, které nebyly součástí výuky (Angular, RxJS), a naučil se pracovat v týmu s využitím pokročilých nástrojů pro správu verzí a CI/CD. Praxe mi potvrdila, že kvalitní teoretický základ ze školy je nezbytný, ale pro uplatnění v oboru je nutné jej neustále rozvíjet samostudiem a reálnou prací na projektech.